\chapter{Result}
\label{ch:results}

\section{Latency and Loop Rates}
\label{sec:results-latency}

\red{Promised in Chapter~\ref{ch:embedded} (System Architecture) and Section~\ref{sec:DAV2}.
Report: loop rate table (PI controller every 10\,ms, encoder count packets every 100\,ms, MDE
range scan and MCL update at about 3\,Hz, navigation loop); DAv2-Small inference speed on the
Raspberry Pi 5 versus the laptop, which justifies offloading depth to the laptop; optionally
the extra $\sim$100\,ms of the metric DAv2 models (Section~\ref{sec:depth-calibration}).}

Loop rate table report

Inference speed of model on embedded software 
thus moving to laptop


\section{Wheel Speed Step Response}
\label{sec:results-step-response}

\red{Promised in Section~\ref{sec:PI}. Report: measured speed $\Delta C[k]$ versus time for a
step from zero to a constant target, for each wheel, with rise time, percent overshoot $M_p$
(Equation~\ref{eq:overshoot}), settling time ($\pm 2\%$), and steady-state error.}

\section{Monocular Depth Accuracy}
\label{sec:results-depth}

\red{Placeholder currently at the end of the Back-Projection to 3D subsection
(Section~\ref{sec:mde-to-range}). Report: per-placement error between checkerboard corners
back-projected from MDE depth and their \texttt{solvePnP} positions, at several distances and
lateral offsets (table or figure).}


\section{Measurement Model Weights}
\label{sec:results-meas-model}

\red{Placeholder currently in Section~\ref{sec:meas-test}. Report: weight maps for each test
case in Table~\ref{tab:meas-tests} (Figure~\ref{fig:meas-test}) and what each case shows about
the sensitivity of the weights to $x$, $y$, and $\theta$. Also the note there about needing a
minimum range and an occupancy value of 0.8.}

\section{Localization Accuracy on the Rover}
\label{sec:results-localization}

\red{Promised in Section~\ref{sec:navigation}. Report: true path extracted from the overhead
tripod camera versus the MCL path (sequence of mean poses), which tests MDE and MCL as one
system. Path comparison figure plus a position and heading error summary.}

\section{Navigation Success Rate}
\label{sec:results-navigation}

\red{Promised in Section~\ref{sec:navigation}. Report: number of successful trials out of
$10$, where a success means the rover stops within $5$~cm of the goal, as measured by the
tripod camera.}

\section{Observed Limitations}
\label{sec:results-limitations}

We found that the minimum range for the particle is a significant drawback of the robot,
since smaller scene give an almost dense map with barely placeable particles. Thus for
diverse tests we have to settle with a more spacious scene. We can mitigate this with the 
abolishing of checkerboard and using pure coordinate transform to get z depth at the closest
points possible but this wont elivate much.


%%% Local Variables: ***
%%% mode: latex ***
%%% TeX-master: "thesis.tex" ***
%%% End: ***
